% !PRINTED 491
\weberpage{491}
%<<CONT
\gesperrt{gen\"ugt; deren Co\"efficienten $A_1, A_2, \ldots, A_m$ ganze Zahlen sind.}

Wir bemerken, dass es nach dieser Definition ausreicht, um eine algebraische Zahl $\Theta$ als ganz zu charakterisiren, wenn unter den unendlich vielen Gleichungen der Form (1), denen $\Theta$ gen\"ugt, \gesperrt{eine} ist, deren Co\"efficienten ganze Zahlen sind.

Die ganzen algebraischen Zahlen umfassen als speciellen Fall die gew\"ohnlichen ganzen Zahlen, die wir zur Unterscheidung \gesperrt{ganze rationale Zahlen} nennen. Die positiven ganzen rationalen Zahlen nennen wir auch, einem verbreiteten Sprachgebrauch folgend, \gesperrt{nat\"urliche Zahlen}.

Unter ganzen Zahlen schlechtweg verstehen wir dann ganze algebraische, rationale und irrationale Zahlen.

1. \gesperrt{Eine ganze algebraische Zahl, die zugleich rational ist, ist nothwendig eine ganze rationale Zahl.}

Nehmen wir n\"amlich an, es sei $\Theta = P:Q$ ein rationaler Bruch, und $P, Q$ ganze rationale Zahlen ohne gemeinsamen Theiler, etwa $Q$ positiv, so ergiebt sich aus (1)
\[
P^m + A_1 P^{m-1} Q + A_2 P^{m-2} Q^2 + \cdots + A_m Q^m = 0,
\]
und daraus ist zu ersehen, dass jeder Primtheiler von $Q$ in $P$ enthalten sein m\"usste. Es muss also $Q = 1$ sein, und $\Theta = P$ ist eine ganze rationale Zahl.

2. \gesperrt{Summe, Differenz und Product zweier ganzer Zahlen sind wieder ganze Zahlen.}

Um diesen Hauptsatz zu beweisen, nehmen wir an, es seien $\alpha, \beta$ zwei ganze Zahlen, die den Gleichungen
\weq{2}{\begin{array}{l}
\alpha^\mu + a_1 \alpha^{\mu-1} + \cdots + a_{\mu-1} \alpha + a_\mu = 0 \\
\beta^\nu + b_1 \beta^{\nu-1} + \cdots + b_{\nu-1} \beta + b_\nu = 0
\end{array}}
gen\"ugen, und machen eine der drei Annahmen
\[
\omega = \alpha + \beta, \ \alpha - \beta, \ \alpha \beta.
\]
Dann setzen wir $\mu \nu = m$ und bezeichnen die $m$ Gr\"ossen
\[
\alpha^r \beta^s \quad \begin{array}{l} r = 0, 1, \ldots, \mu - 1 \\ s = 0, 1, \ldots, \nu - 1 \end{array}
\]
in irgend einer Reihenfolge mit $\omega_1, \omega_2, \ldots, \omega_m$.
